\documentclass[11pt,a4paper]{article}

% --- PDFLATEX COMPATIBLE PREAMBLE ---
\usepackage[utf8]{inputenc}
\usepackage[T1]{fontenc}
\usepackage[english]{babel}
\usepackage{helvet}         
\renewcommand{\familydefault}{\sfdefault}

% Adjusted margins to fit everything nicely on one page
\usepackage[left=0.75in,top=0.6in,right=0.75in,bottom=0.6in]{geometry}
\usepackage{url}
\usepackage{hyperref}
\usepackage{enumitem}
\usepackage{xcolor}

% Custom Colors 
\definecolor{TECHBLUE}{RGB}{0, 0, 0} 

% Prevent hyphenation 
\hyphenpenalty=10000
\exhyphenpenalty=10000
\sloppy

\hypersetup{
    colorlinks=true,
    linkcolor=black,
    urlcolor=black,
}

\pagestyle{empty}
\setlength{\parindent}{0pt}
\setlength{\parskip}{0.5em} % Reduced slightly to fit one page

\begin{document}
%-----------------------------------------------------------
% Sender Information
%-----------------------------------------------------------
\textbf{\large Juan David Muñoz-Bolaños} \\
Rechenauerstraße 42, Innsbruck, Austria (Tirol) \\
+43 676 9115145 \\
\href{mailto:juan.munoz@i-med.ac.at}{juan.munoz@i-med.ac.at}

\vspace{8pt}

%-----------------------------------------------------------
% Recipient Information
%-----------------------------------------------------------
\textbf{FERCHAU GmbH} \\
Frau Laura Längsfeld \\
Niederlassung Rosenheim \\
Deutschland \\
 \\


\vspace{8pt}

%-----------------------------------------------------------
% Subject Line
%-----------------------------------------------------------
\textbf{Subject: Bewerbung als Licht-/Optikentwickler (m/w/d) – Kennziffer FE37-48083-RH}

\vspace{4pt}

%-----------------------------------------------------------
% Body
%-----------------------------------------------------------
Sehr geehrte Frau Längsfeld,


Grosse Begeisterung für opto-mechanisches Design und photonische Simulation treiben mich an, mich für die Position als Licht- und Optikentwickler bei FERCHAU in Glonn zu bewerben. Die End-to-End-Entwicklung optischer Konzepte – von der Simulation bis zum Laborversuch – entspricht genau meiner Expertise.



In meiner Promotion an der Medizinischen Universität Innsbruck habe ich komplexe optische Systeme inklusive Beleuchtung und Laserintegration entwickelt und qualifiziert. Dabei gehörten Design, Aufbau, Validierung von lichttechnischen Konzepten und Toleranzmanagement eng zusammen. Zuvor am Leibniz-Institut für Photonik realisierte ich faseroptische Sonden. In beiden Stationen koordinierte ich Auswahl und Beschaffung intensiv mit Lieferanten.



Aus meiner Zeit als Machine Vision Engineer bei INTECOL S.A.S. bringe ich zudem Industrieerfahrung mit. Dort verantwortete ich massgeschneiderte Beleuchtungssysteme (LED/Ringlichter/Dunkelfeld) für optische Qualitätskontrollen sowie den Aufbau normgerechter Laborversuche. Analytisches Denken und fundierte MS Office-Dokumentation sind für mich selbstverständlich.



Ursprünglich aus Kolumbien, habe ich meinen Master in Jena absolviert und promoviere nun in Innsbruck, wodurch ich im DACH-Raum bestens integriert bin. An meinem Deutsch (B2) arbeite ich täglich mit Engagement. Die Nähe zu den Bergen schätze ich sehr, weshalb mich der Standort Glonn besonders reizt.



Mit Promotionsabschluss im Mai bin ich kurzfristig verfügbar, reisebereit und besitze einen B-Führerschein. Ich freue mich, Ihr Team mit interdisziplinärem Know-how zu verstärken.




\vspace{12pt}

%-----------------------------------------------------------
% Closing
%-----------------------------------------------------------
Mit freundlichen Grüssen,

\vspace{8pt}

\textbf{Juan David Muñoz-Bolaños} \\
Licht- und Optikentwickler \\
Medizinische Universität Innsbruck

\end{document}